\scanpage{3}
% slanted
\margin{\uline{Vanishing} cycles : cas des quadriques}

\diagram{\begin{tikzcd}[ampersand replacement=\&]
 X \arrow[d,"f"] \\ S
\end{tikzcd}}

$f$ un schéma projectif plat

$S$ spectre d'un anneau de val. discrète strictement local

$X_0$ fibre spéciale quadrique à point singulier

$\bar X_1$ fibre géom générale, quadrique non singulière

dimension (commune) des fibres $n-1$

\unsure{$\mathrm{I}$|$[$}

\deleted{On veut étudier}

$X\deleted{\unsure{=|{-}}}\subset P^n_S$ défini par une équation quadratique \ldots.

$H^i(X_0)\to H^i(\bar X_1)$ ? On a la suite exacte générale

\[
\cdots\to H^i(X_0)\to H^i(\bar X_1)\unsure{\mathord{-}|\to}\Phi^{i+1}\to H^{i+1}(X_0)\to H^{i+1}(\bar X_1)\to\cdots
\]

\deleted{1\textsuperscript{o}). \uline{$n$ pair}, d'où \unsure{$n-1$ impair|$n$ impair}}

On sait que la cohomologie de \deleted{$H^i(X_0)$ et} $X_0$ et $\bar X_1$ est nulle en dims impaires. Donc on trouve

\[
0\to\Phi^{2\nu}\to H^{2\nu}(X_0)\to H^{2\nu}(\bar X_1)\to\Phi^{2\nu+1}\to0
\]

i.e. les $\Phi^{i}$ apparaissent \deleted{\unsure{certes|comme}} comme les noyaux et conoyaux des \uline{hom de spécialisation}.

Notons aussi que sur $X$ donc $X_0$, on a une classe $\xi\in H^2(X,\mu(1))$, provenant de la section à l'infini, et dont la restriction aux deux fibres est connue \ldots

% archive sheet number: 2
